% ======================================================================
% SECTION 1: HR 9501 - Cancer Patient Self-Selection of Physical AI
% Oncology Trials Act of 2026
% Adaption of 21 CFR Part 50, FDA Decentralized Trial Final Guidance
% 2024, and 42 USC 300gg-8.
%
% Layout follows the prior site documentation template at
% new-trial/site/all-documents/all_documents.tex (Findings,
% Definitions, Authorization or Rights, Implementation, Reporting and
% Enforcement). The section title in main.tex names the bill as an
% Adaption.
%
% This file consolidates instructions previously housed in the prior
% template's Abstract, Introduction, Patient Priority Framework
% (dimensions 2.1 trial discovery and 2.3 location/scheduling),
% Implementation/Metrics (Domain 1 access, timeline tie-in),
% Discussion (11.1.b CLINICAL TREATMENT Act extension; 11.1.c ACA
% Section 2709 extension; 11.2.a LLM trial-matching evidence),
% Limitations (12.1.a HR 4501 limitation paragraph), and Conclusions
% (13.1 first list item) into the body of HR 9501.
% ======================================================================

[CONSOLIDATED OPENING THESIS (formerly Abstract opening sentence,
moved here as the lead bill in the paper). The cancer patient, NOT the
doctor, the nurse, the trial sponsor, the IRB, or the regulator, is
the priority participant in any U.S. oncology clinical trial. Advanced
AI and advanced robotics now offer the patient direct control over
discovery, enrollment, scheduling, robot selection, procedural
modification, real-time data feedback, and self-custody of health
data. HR 9501 is the foundational bill of seven proposed federal bills
(HR 9501 through HR 9507) that supply the legal infrastructure for
that control. Cite \cite{kawchak_2026_19994945},
\cite{kawchak_2026_19119939}, and \cite{kawchak_2026_20045457}.]

[CONSOLIDATED CURRENT-LEGISLATION CONTEXT (formerly Introduction
block 2 items 2.b, 2.c, 2.e, 2.h relevant to trial discovery and
self-selection). Read in full
patients/paper/Deep-Research-2/chunk_01_intro_ranks1-4.md (Ranks 1-4:
21st Century Cures Act PL 114-255; CLINICAL TREATMENT Act PL 116-260;
ACA Section 2709 PL 111-148; FDORA / FDA Modernization Act 2.0 PL
117-328 Section 3209) and
patients/paper/Deep-Research-2/chunk_02_ranks5-9.md (Ranks 5-9: FDAAA
PL 110-85 Section 801; Right to Try Act PL 115-176; FDARA PL 115-52;
CA SB 37 (2001); ONC HTI-1 Final Rule 2023). Synthesize the
patient-control deficits HR 9501 closes:

(1.0.a) CLINICAL TREATMENT Act, PL 116-260 (2020): Medicaid coverage
of routine costs in qualifying trials from January 1 2022. HR 9501
extends this coverage right into a positive self-selection right that
includes 24/7 booking, decentralized payment parity, and AI trial-
matching tool actionability for Medicaid-eligible patients. Cite
\cite{rank2_clinical_treatment_act_2020}, \cite{jin2024trialgpt},
\cite{gupta2024prism}, and \cite{abdallah2026trialmatchai}.

(1.0.b) Affordable Care Act Section 2709, PL 111-148 (2010): private
insurance baseline for routine clinical-trial cost coverage. HR 9501
closes the decentralized and at-home component payment-parity gap
that ACA Section 2709 left unaddressed. Cite
\cite{rank3_aca_clinical_trials_2010} and
\cite{usc_300gg8_clinical_trials}.

(1.0.c) FDAAA Section 801, PL 110-85 (2007): the ClinicalTrials.gov
infrastructure. HR 9501 (jointly with HR 9507) closes the structured-
eligibility-data gap that prevents patient-initiated trial discovery.
Cite \cite{rank5_fdaaa_2007}.

(1.0.d) HHS HIPAA Right-to-Access (2025), ASTP/ONC Patient Portals
Brief 2024, and ASTP/ONC Hospital APIs Brief 2024. Reference these
foundational frameworks but defer the self-custody right to HR 9507.
Cite \cite{hhs_right_to_access_2025},
\cite{astp_onc_patient_portals_apps_2024}, and
\cite{astp_onc_hospital_apis_2024}.]

[CONSOLIDATED PATIENT PRIORITY FRAMEWORK DIMENSIONS (formerly Patient
Priority Framework subsections 2.1 and 2.3). HR 9501 implements two
of the seven dimensions of the operational definition of patient
control:

Dimension 2.1 - Trial discovery: machine-readable eligibility, AI
matching, no provider gatekeeping. Cite \cite{jin2024trialgpt},
\cite{gupta2024prism}, and \cite{abdallah2026trialmatchai}.

Dimension 2.3 - Location and scheduling choice: 24/7 booking,
decentralized elements, home-based delivery. Cite
\cite{fda_dct_guidance_2024},
\cite{fda_oce_advancing_oncology_dct_2024},
\cite{dronca2026cancer_care_beyond_walls}, and
\cite{kawchak_2026_19994945} (for the prior 24/7 trial paper).]

[PROCESSING INSTRUCTION (HR 9501 - Cancer Patient Self-Selection of
Physical AI Oncology Trials Act of 2026, populate to 9 to 12
paragraphs structured into the five legislative-act subsections):
Draft the proposed bill text in legislative-act voice, following the
layout used by SB 1042 in the prior site template at
new-trial/site/all-documents/all_documents.tex (which was, in turn,
populated from
new-trial/site/01-legislation-authorization/physical_ai_trial_authorization.tex
and national-platform/new_trial_psl/01_preamble_and_sb1042.tex).

The bill must explicitly state in its first paragraph that it is an
"Adaption of 21 CFR Part 50, FDA Decentralized Trial Final Guidance
2024, and 42 USC 300gg-8" so that readers can trace the chain of prior
authority. Cite \cite{hhs_45cfr46_protection_human_subjects},
\cite{fda_dct_guidance_2024}, \cite{fda_oce_advancing_oncology_dct_2024},
and \cite{usc_300gg8_clinical_trials}.

Subsection 1.1 - Findings and Declarations: Read
patients/paper/Deep-Research-3/part1_legal_baseline.md and quote the
seven percent participation-rate finding. Read
patients/paper/Deep-Research-3/part2_ai_robotics_legislation.md and
quote the TrialGPT, PRISM, and TrialMatchAI performance figures.
Compose 5 to 7 enumerated findings in the (a)/(b)/(c) style used by
SB 1042 in the prior template. Each finding must cite at least one
prior-law reference and at least one AI-evidence reference. Cite
\cite{jin2024trialgpt}, \cite{gupta2024prism}, and
\cite{abdallah2026trialmatchai}.

Subsection 1.2 - Definitions: Define the following operational terms
used throughout the bill:
- "Patient self-selection" means the patient's right to discover,
  evaluate, and book their own oncology clinical trial slot without
  intermediary provider gatekeeping. Reference the existing 24/7 trial
  format documented in
  new-trial/national-24-7-trial/paper/full-paper/main.tex and the
  hour-NN/patient_records.md files at
  new-trial/national-24-7-trial/hour-00/ through hour-55/.
- "Machine-readable eligibility report" means a structured FHIR-based
  output containing the patient's eligibility-relevant data fields
  for at least 80 percent of currently recruiting U.S. oncology
  trials. Reference the ClinicalTrials.gov infrastructure under FDAAA
  Section 801. Cite \cite{rank5_fdaaa_2007}.
- "Qualified Physical AI oncology trial site" means a site authorized
  under either this Act or under the prior California SB 1042 / Title
  22 framework documented at
  new-trial/site/05-state-regulations/ca_state_regulations.tex and
  national-platform/new_trial_psl/05_title22_regulations.tex (read
  this file in full when populating).
- "24/7 booking" means the patient's right to book any qualified site
  at any hour of the day, any day of the week, drawing on the
  168-hour sponsor extension at sponsor/final_paper/168_hours/ for
  the operational baseline. Cite \cite{kawchak_2026_19396256} and
  \cite{kawchak_2026_19994945}.

Subsection 1.3 - Patient Self-Selection Rights (the operative core):
Compose 6 to 8 enumerated rights in the (a)/(b)/(c) style. Each right
must be specific enough to be enforceable. Examples:
- Right to receive a machine-readable eligibility report from any
  CMS-billable EHR within 7 calendar days of pathology or genomic
  result release.
- Right to discover all currently recruiting U.S. oncology trials via
  an FDA-published API that integrates the data fields from
  ClinicalTrials.gov, the FDA RTCT pilot program, and FDA Oncology
  Center of Excellence DCT directory at
  new-trial/national-24-7-trial/FDA-April-2026/FDA_RealTime_Clinical_Trials.md
  and \cite{fda_oce_advancing_oncology_dct_2024}.
- Right to book a 24/7 trial slot at any qualified site without prior
  written referral. Reference the 24/7 site infrastructure at
  new-trial/national-24-7-trial/hour-NN/diagram_facility.txt files
  (e.g., hour-00, hour-12, hour-23, hour-47) for the operational
  pattern.
- Right to payment parity for decentralized and at-home trial
  components under 42 USC 300gg-8 and CMS NCD 310.1; closes the gap
  identified in the consolidated current-legislation context above.
  Cite \cite{usc_300gg8_clinical_trials} and
  \cite{cms_ncd_3101_routine_costs}.
- Right to cross-state portability of eligibility data and trial
  bookings, drawing on the ASTP/ONC Hospital APIs Brief 2024 to
  identify the technical baseline. Cite
  \cite{astp_onc_hospital_apis_2024}.
- Right to refuse provider-only routing of trial discovery in favor
  of an AI matching tool of the patient's choosing (TrialGPT, PRISM,
  TrialMatchAI), with disclosure of model provenance under HTI-1 DSI
  transparency. Cite \cite{hti1_dsi_fact_sheet_2023}.

Subsection 1.4 - Implementation: 12-month deadline for the federal
API; 24-month deadline for full state-level adoption; FDA, ONC, HHS,
and CMS roles; preemption rule that explicitly preserves stronger
state-level patient protections (e.g., CA SB 37 for routine-cost
coverage). Cite \cite{rank8_california_sb37_2001} and
\cite{rank4_fdora_fda_modernization_2022} for the FDORA precedent
that authorizes mandatory FDA decentralized-trial guidance.

Subsection 1.5 - Reporting and Enforcement: Annual public report on
patient-self-selection rates by state, payer, and cancer type;
information-blocking penalties harmonized with
\cite{federal_register_info_blocking_disincentives_2024}; an explicit
"no automated-only denial" clause that mirrors the first guardrail
from patients/paper/Deep-Research-3/part3_metrics_guardrails.md.]

[CONSOLIDATED IMPLEMENTATION-METRICS DOMAIN 1 - ACCESS (formerly
Implementation Metrics subsection 10.2 Domain 1, moved here because
HR 9501 is the access bill). The federal Patient Control Dashboard
must publish, for HR 9501, the following Domain 1 metrics:
trial-offer rate, days from pathology or genomic result to trial
shortlist, machine-readable eligibility report coverage, real-time
data release rates, travel and work hours saved. Each metric is
drawn from
patients/paper/Deep-Research-3/part3_metrics_guardrails.md. Cite
\cite{hhs_right_to_access_2025} and
\cite{rocque2025remote_symptom_monitoring}.]

[CONSOLIDATED DISCUSSION COMPARISON (formerly Discussion items 11.1.b
CLINICAL TREATMENT Act extension and 11.1.c ACA Section 2709
extension and 11.2.a LLM trial-matching evidence). HR 9501 extends
the prior CLINICAL TREATMENT Act Medicaid coverage of routine costs
into a positive self-selection right that includes 24/7 booking and
decentralized payment parity, closing the gap that ACA Section 2709
left unaddressed. The LLM trial-matching evidence (TrialGPT, PRISM,
TrialMatchAI) demonstrates the technology is ready, but HR 9501
(self-selection) and HR 9507 (data self-custody) are required before
patients can use it. Cite \cite{rank2_clinical_treatment_act_2020},
\cite{rank3_aca_clinical_trials_2010},
\cite{cms_ncd_3101_routine_costs}, \cite{jin2024trialgpt},
\cite{gupta2024prism}, and \cite{abdallah2026trialmatchai}.]

[CONSOLIDATED LIMITATIONS PARAGRAPH (formerly Limitations item
12.1.a). HR 9501 assumes machine-readable eligibility data of high
quality across all CMS-billable EHRs. The ASTP/ONC hospital APIs
brief documents that adoption is uneven; the 12-month deadline is
aggressive. The Mazor 2025 null result on AI notifications alone
reminds us that self-selection rights without a complete pathway
from discovery to consent to enrollment may not improve participation
on their own. Cite \cite{astp_onc_hospital_apis_2024} and
\cite{mazor2025ai_notifications_trial}.]

[CONSOLIDATED CONCLUSIONS RESTATEMENT (formerly Conclusions item
13.1 first list item). HR 9501 closes the patient-control gap left
by the layered legal stack (Cures Act, CLINICAL TREATMENT Act, ACA
Section 2709, FDAAA Section 801) by replacing provider-gatekept trial
referral with patient-initiated, AI-matched, 24/7-bookable
self-selection across all qualified Physical AI oncology trial sites.
Cite \cite{rank1_cures_act_2016},
\cite{rank2_clinical_treatment_act_2020}, and
\cite{kawchak_2026_19994945}.]

Close the bill (1 paragraph after Subsection 1.5) with two patient-
control framings:

(i) FOR INDIVIDUAL PATIENTS: An NSCLC patient at home in rural Texas
    can, on their own initiative, receive an FHIR eligibility report
    for the qualified Physical AI Oncology Trial Site documented in
    national-platform/new_trial_psl/01_preamble_and_sb1042.tex within
    7 days of their genomic result, then self-book a 24/7 slot at a
    USL-7-or-above surgical robot system following the 10-stage
    journey for PAT-2026-0042 in
    patient-journey/stage_01_prescreening.py through
    stage_10_closeout.py.

(ii) FOR BROAD ADOPTION: All U.S. oncology trials must publish their
     eligibility criteria in machine-readable form within the
     12-month federal-API window, with ALL CMS-billable EHRs
     maintaining a self-selection portal endpoint accessible by every
     U.S. cancer patient by month 24. The United States must remain
     Number 1 in the world for cancer patient benefit during the
     medical AI and robotics revolution, and HR 9501 is the
     foundational bill that opens that pathway. Cite
     \cite{kawchak_2026_19994945} and \cite{kawchak_2026_19119939}.
